\lecture{4.6}{G-Computation and IPTW in R}
\label{lec:4-6}

\subsection{Data Preparation}

The R implementations of g-computation and IPTW are most reliable when the data are
prepared so that the notation matches the analysis. The dataset should contain one row
per individual and variables corresponding to $W$, $A$, and $Y$:
\[
O_i=(W_i,A_i,Y_i).
\]
Treatment should be coded consistently, often as $A=1$ for treated and $A=0$ for
untreated. Binary outcomes should also be coded as $0/1$ when logistic regression is
used.

\begin{implementationbox}
Example setup:
\begin{lstlisting}[language=R]
# dat has one row per individual
# A: binary treatment, coded 0/1
# Y: outcome
# W1, W2: pre-treatment covariates

dat <- dat[complete.cases(dat[, c("Y", "A", "W1", "W2")]), ]
dat$A <- as.numeric(dat$A == 1)
dat$Y <- as.numeric(dat$Y)
\end{lstlisting}
\end{implementationbox}

Dropping observations with missing data is simple but not always appropriate. If
missingness is related to treatment, outcome, or covariates, complete-case analysis can
change the target population or introduce bias. Later chapters treat missing data more
carefully.

\subsection{Fitting Outcome Regressions}

For g-computation, the nuisance function is
\[
\Qbar(a,w)=\E(Y\given A=a,W=w).
\]
The fitted outcome regression should include treatment and the confounders required
for identification. It may also include interactions or nonlinear terms if scientifically
or statistically warranted.

\begin{implementationbox}
Outcome regression examples:
\begin{lstlisting}[language=R]
# Continuous outcome, identity link
fit_Q_linear <- glm(Y ~ A + W1 + W2, data = dat, family = gaussian())

# Binary outcome, logistic link
fit_Q_logistic <- glm(Y ~ A + W1 + W2, data = dat, family = binomial())

# Interaction model allowing treatment effect heterogeneity by W1
fit_Q_interact <- glm(Y ~ A * W1 + W2, data = dat, family = binomial())
\end{lstlisting}
\end{implementationbox}

The fitted object is not itself the causal estimate. It is an estimate of a nuisance
function used to compute predicted counterfactual outcomes.

\subsection{Fitting Propensity Score Models}

For IPTW, the nuisance function is
\[
g_1(w)=\Prob(A=1\given W=w).
\]
Because $A$ is binary, logistic regression is a common starting point:
\[
\logit\{g_1(w)\}=\alpha_0+\alpha^\top w.
\]

\begin{implementationbox}
Propensity score model:
\begin{lstlisting}[language=R]
fit_g <- glm(A ~ W1 + W2, data = dat, family = binomial())

g1 <- predict(fit_g, type = "response")
g0 <- 1 - g1
\end{lstlisting}
\end{implementationbox}

After fitting the propensity score, inspect whether predicted probabilities are close
to zero or one. Very small $g_1(W_i)$ among treated individuals or very small
$g_0(W_i)$ among untreated individuals lead to large IPTW weights.

\subsection{Prediction Workflows}

The g-computation prediction workflow constructs the two estimated counterfactual
means by changing only the treatment variable:

\[
\psi_n(1)=\frac{1}{n}\sum_{i=1}^n\Qbar_n(1,W_i),
\qquad
\psi_n(0)=\frac{1}{n}\sum_{i=1}^n\Qbar_n(0,W_i).
\]

\begin{implementationbox}
G-computation prediction workflow:
\begin{lstlisting}[language=R]
dat1 <- dat
dat1$A <- 1

dat0 <- dat
dat0$A <- 0

Q1 <- predict(fit_Q_logistic, newdata = dat1, type = "response")
Q0 <- predict(fit_Q_logistic, newdata = dat0, type = "response")

psi1_gcomp <- mean(Q1)
psi0_gcomp <- mean(Q0)
ate_gcomp  <- psi1_gcomp - psi0_gcomp
\end{lstlisting}
\end{implementationbox}

For IPTW, no treatment-specific prediction datasets are required. The estimator uses
the observed treatment and observed outcome:
\[
\psi_n^{\IPTW}(a)=
\frac{1}{n}\sum_{i=1}^n
\frac{\ind(A_i=a)}{g_{n,a}(W_i)}Y_i.
\]

\begin{implementationbox}
IPTW workflow:
\begin{lstlisting}[language=R]
A <- dat$A
Y <- dat$Y

psi1_iptw <- mean((A == 1) / g1 * Y)
psi0_iptw <- mean((A == 0) / g0 * Y)
ate_iptw  <- psi1_iptw - psi0_iptw
\end{lstlisting}
\end{implementationbox}

\subsection{Treatment-Specific Datasets}

Treatment-specific datasets are useful because they make the estimand visible in the
code. The dataset \texttt{dat1} represents the empirical covariate distribution with
everyone assigned $A=1$; \texttt{dat0} represents the same empirical covariate
distribution with everyone assigned $A=0$.

\begin{figure}[H]
\centering
\begin{tikzpicture}[
  node distance=1.35cm,
  box/.style={draw, rounded corners=2pt, minimum width=3.2cm, minimum height=.8cm},
  arr/.style={-{Latex[length=2.5mm]}, thick}
]
\node[box] (dat) {Observed covariates\\ $W_1,\ldots,W_n$};
\node[box, below left=1cm and .9cm of dat] (dat1) {Set all\\ $A=1$};
\node[box, below right=1cm and .9cm of dat] (dat0) {Set all\\ $A=0$};
\node[box, below=1cm of dat1] (q1) {Predict\\ $\Qbar_n(1,W_i)$};
\node[box, below=1cm of dat0] (q0) {Predict\\ $\Qbar_n(0,W_i)$};
\node[box, below right=1cm and .15cm of q1] (ate) {Average\\ and subtract};
\draw[arr] (dat) -- (dat1);
\draw[arr] (dat) -- (dat0);
\draw[arr] (dat1) -- (q1);
\draw[arr] (dat0) -- (q0);
\draw[arr] (q1) -- (ate);
\draw[arr] (q0) -- (ate);
\end{tikzpicture}
\caption{Treatment-specific datasets implement g-computation predictions.}
\label{fig:ch4-treatment-specific-data}
\end{figure}

This workflow should not be confused with changing the observed data. The original
outcomes remain the data used to fit the outcome model. The modified datasets are
only prediction grids used to evaluate the fitted conditional mean.

\subsection{Practical Implementation Patterns}

It is often helpful to wrap estimators in functions. This makes bootstrap inference
straightforward and reduces the chance of accidentally using variables from the wrong
dataset.

\begin{implementationbox}
Reusable g-computation function:
\begin{lstlisting}[language=R]
get_gcomp <- function(input_data) {
  fit_Q <- glm(Y ~ A + W1 + W2, data = input_data, family = binomial())

  data1 <- input_data
  data1$A <- 1
  data0 <- input_data
  data0$A <- 0

  Q1 <- predict(fit_Q, newdata = data1, type = "response")
  Q0 <- predict(fit_Q, newdata = data0, type = "response")

  c(psi0 = mean(Q0), psi1 = mean(Q1), ate = mean(Q1) - mean(Q0))
}
\end{lstlisting}
\end{implementationbox}

\begin{implementationbox}
Reusable IPTW function:
\begin{lstlisting}[language=R]
get_iptw <- function(input_data) {
  fit_g <- glm(A ~ W1 + W2, data = input_data, family = binomial())

  g1 <- predict(fit_g, type = "response")
  g0 <- 1 - g1

  A <- input_data$A
  Y <- input_data$Y

  psi1 <- mean((A == 1) / g1 * Y)
  psi0 <- mean((A == 0) / g0 * Y)

  c(psi0 = psi0, psi1 = psi1, ate = psi1 - psi0)
}
\end{lstlisting}
\end{implementationbox}

\begin{estimationbox}
Implementation checks:
\begin{itemize}
\item Confirm that $A$, $Y$, and all covariates are coded as intended.
\item For g-computation, confirm that \texttt{newdata} contains the same covariate
names and factor levels as the model-fitting dataset.
\item For IPTW, inspect summaries of $g_1$, $g_0$, and the resulting weights.
\item Keep the estimation function self-contained so that resampling methods refit
all nuisance functions inside each resampled dataset.
\end{itemize}
\end{estimationbox}
